\inhaltkopf{INHALT DES ZWEITEN BANDES.}

\inhaltbuch{Erstes Buch.}{Gruppen.}

\inhaltabschnitt{Erster Abschnitt.}{Allgemeine Gruppentheorie.}
\inhaltpar{1}{Definition der Gruppen}{3}
\inhaltpar{2}{Die Divisoren endlicher Gruppen}{7}
\inhaltpar{3}{Normaltheiler einer Gruppe}{10}
\inhaltpar{4}{Composition der Theile}{12}
\inhaltpar{5}{Mehrstufiger Isomorphismus}{15}
\inhaltpar{6}{Die Compositionsreihe und der Satz von C. \gesperrt{Jordan}}{17}
\inhaltpar{7}{Weitere Sätze über die Compositionsreihen}{24}
\inhaltpar{8}{Metacyklische Gruppen}{27}

\inhaltabschnitt{Zweiter Abschnitt.}{Abel'sche Gruppen.}
\inhaltpar{9}{Darstellung \gesperrt{Abel}'scher Gruppen durch eine Basis}{32}
\inhaltpar{10}{Die Invarianten der \gesperrt{Abel}'schen Gruppen}{39}
\inhaltpar{11}{Gruppencharaktere}{43}
\inhaltpar{12}{Divisoren einer \gesperrt{Abel}'schen Gruppe. Reciproke Gruppen}{48}
\inhaltpar{13}{Die Geschlechter in einer \gesperrt{Abel}'schen Gruppe}{52}
\inhaltpar{14}{Indices nach einer ungeraden Primzahlpotenz als Modul}{54}
\inhaltpar{15}{Indices für eine Potenz von 2 als Modul}{58}
\inhaltpar{16}{Die Gruppen der Zahlclassen nach einem zusammengesetzten Modul}{60}

\inhaltabschnitt{Dritter Abschnitt.}{Die Gruppe der Kreistheilungskörper.}
\inhaltpar{17}{Die Resolventen der Kreistheilungstheorie}{63}
\inhaltpar{18}{Kreistheilungskörper}{67}
\inhaltpar{19}{Primäre und nicht primäre Theiler der Gruppe $\mathfrak{N}$}{72}
\inhaltpar{20}{Die Kreistheilungsperioden}{74}
\inhaltpar{21}{Kreistheilungskörper mit gegebener Gruppe}{79}
\inhaltpar{22}{Bestimmung der Gruppe $\mathfrak{A}$}{92}
